\documentclass[10pt,a4paper]{amsart}
\usepackage{amsmath,amssymb,amsthm,graphicx}
% Retained minimal preamble for the Stacks Project chapter more-morphisms.tex (More on Morphisms).
% Verbatim excerpt of the chapter preamble: Xy-pic setup (native preamble lines 14-18) and the native
% theorem declarations with the mathematical macros (native preamble lines 154-217). Copied mechanically;
% nothing added, removed or redefined. The wrapper class and the amsmath/amssymb/amsthm/graphicx base
% packages are supplied by the assembler, so the class/externaldocument part of the native preamble is omitted.
% For commutative diagrams we use Xy-pic
\usepackage[all]{xy}

% We use 2cell for 2-commutative diagrams.
\xyoption{2cell}
% Theorem environments.
%
\theoremstyle{plain}
\newtheorem{theorem}[subsection]{Theorem}
\newtheorem{proposition}[subsection]{Proposition}
\newtheorem{lemma}[subsection]{Lemma}

\theoremstyle{definition}
\newtheorem{definition}[subsection]{Definition}
\newtheorem{example}[subsection]{Example}
\newtheorem{exercise}[subsection]{Exercise}
\newtheorem{situation}[subsection]{Situation}

\theoremstyle{remark}
\newtheorem{remark}[subsection]{Remark}
\newtheorem{remarks}[subsection]{Remarks}

\numberwithin{equation}{subsection}

% Macros
%
\def\lim{\mathop{\mathrm{lim}}\nolimits}
\def\colim{\mathop{\mathrm{colim}}\nolimits}
\def\Spec{\mathop{\mathrm{Spec}}}
\def\Hom{\mathop{\mathrm{Hom}}\nolimits}
\def\Ext{\mathop{\mathrm{Ext}}\nolimits}
\def\SheafHom{\mathop{\mathcal{H}\!\mathit{om}}\nolimits}
\def\SheafExt{\mathop{\mathcal{E}\!\mathit{xt}}\nolimits}
\def\Sch{\mathit{Sch}}
\def\Mor{\mathop{\mathrm{Mor}}\nolimits}
\def\Ob{\mathop{\mathrm{Ob}}\nolimits}
\def\Sh{\mathop{\mathit{Sh}}\nolimits}
\def\NL{\mathop{N\!L}\nolimits}
\def\CH{\mathop{\mathrm{CH}}\nolimits}
\def\proetale{{pro\text{-}\acute{e}tale}}
\def\etale{{\acute{e}tale}}
\def\QCoh{\mathit{QCoh}}
\def\Ker{\mathop{\mathrm{Ker}}}
\def\Im{\mathop{\mathrm{Im}}}
\def\Coker{\mathop{\mathrm{Coker}}}
\def\Coim{\mathop{\mathrm{Coim}}}

% Boxtimes
%
\DeclareMathSymbol{\boxtimes}{\mathbin}{AMSa}{"02}

%
% Macros for moduli stacks/spaces
%
\def\QCohstack{\mathcal{QC}\!\mathit{oh}}
\def\Cohstack{\mathcal{C}\!\mathit{oh}}
\def\Spacesstack{\mathcal{S}\!\mathit{paces}}
\def\Quotfunctor{\mathrm{Quot}}
\def\Hilbfunctor{\mathrm{Hilb}}
\def\Curvesstack{\mathcal{C}\!\mathit{urves}}
\def\Polarizedstack{\mathcal{P}\!\mathit{olarized}}
\def\Complexesstack{\mathcal{C}\!\mathit{omplexes}}
% \Pic is the operator that assigns to X its picard group, usage \Pic(X)
% \Picardstack_{X/B} denotes the Picard stack of X over B
% \Picardfunctor_{X/B} denotes the Picard functor of X over B
\def\Pic{\mathop{\mathrm{Pic}}\nolimits}
\def\Picardstack{\mathcal{P}\!\mathit{ic}}
\def\Picardfunctor{\mathrm{Pic}}
\def\Deformationcategory{\mathcal{D}\!\mathit{ef}}
\usepackage{cleveref}
\crefname{originallemma}{original lemma}{original lemmas}
\Crefname{originallemma}{Original lemma}{Original lemmas}
\makeatletter
\newlabel{lemma-decompose}{{0GTL}{source}{}{sourcepointer-1}{}}
\newlabel{lemma-decompose@cref}{{[originallemma][0][]0GTL}{[1][0][]source}{}{}{}}
\makeatother
\crefname{originalreference}{original reference}{original references}
\Crefname{originalreference}{Original reference}{Original references}
\makeatletter
\newlabel{limits-lemma-descend-finite-presentation}{{01ZM}{source}{}{sourcepointer-2}{}}
\newlabel{limits-lemma-descend-finite-presentation@cref}{{[originalreference][0][]01ZM}{[1][0][]source}{}{}{}}
\makeatother
\crefname{originalreference}{original reference}{original references}
\Crefname{originalreference}{Original reference}{Original references}
\makeatletter
\newlabel{limits-lemma-descend-immersion}{{0GTB}{source}{}{sourcepointer-3}{}}
\newlabel{limits-lemma-descend-immersion@cref}{{[originalreference][0][]0GTB}{[1][0][]source}{}{}{}}
\makeatother
\crefname{originalreference}{original reference}{original references}
\Crefname{originalreference}{Original reference}{Original references}
\makeatletter
\newlabel{limits-lemma-descend-surjective}{{07RR}{source}{}{sourcepointer-4}{}}
\newlabel{limits-lemma-descend-surjective@cref}{{[originalreference][0][]07RR}{[1][0][]source}{}{}{}}
\makeatother
\crefname{originalreference}{original reference}{original references}
\Crefname{originalreference}{Original reference}{Original references}
\makeatletter
\newlabel{limits-proposition-characterize-locally-finite-presentation}{{01ZC}{source}{}{sourcepointer-5}{}}
\newlabel{limits-proposition-characterize-locally-finite-presentation@cref}{{[originalreference][0][]01ZC}{[1][0][]source}{}{}{}}
\makeatother
\title{Descent of the completely decomposed property along a limit of schemes with affine transitions}
\author{The Stacks Project Authors}
\date{stacks-project:more-morphisms.tex:local-snapshot-20261011; GNU FDL 1.2 or later}
\begin{document}
\maketitle
\noindent\textit{Editorial note.} Counted span is the author's own proof environment, source lines 23257-23290, immediately following the statement at 23239-23255.\par\medskip
\noindent\textit{Editorial note.} Four cross-chapter pointers (Limits) plus a same-chapter lemma are recorded in assembly.external\_references.\par\medskip
\noindent\textit{Editorial note.} The statement fixes the indexing convention (0 in Lambda, base change to S\_lambda); that wording is retained verbatim in the statement span.\par\medskip
\noindent\textit{Allowed standard prerequisites (editorial).}\begin{itemize}
\item Decomposition of completely decomposed morphisms. chapter lemma lemma-decompose, cited in the proof. Provides the decomposition criterion that the descended data has to satisfy.
\item Local finite presentation along a limit. Limits, proposition limits-proposition-characterize-locally-finite-presentation, cited in the proof. Certifies local finite presentation from the presentation at a stage.
\item Descending finite presentation, immersions and surjectivity. Limits, lemmas limits-lemma-descend-finite-presentation, limits-lemma-descend-immersion and limits-lemma-descend-surjective, cited in the proof. Move the finitely presented pieces, the immersions and the surjection from X to some stage X\_lambda.
\end{itemize}\medskip
\ifdefined\hypertarget\hypertarget{sourcepointer-1}{}\fi
\noindent\textit{Source reference.} Original lemma 0GTL: sources/stacks-more-morphisms/pinned-original-source.tex, source lines 23187--23187. Navigation pointer only: the counted proof cites this statement by \textbackslash{}ref as a standing input; the deduction the counted proof has to make is carried out inside the counted span, and no part of that deduction is delegated to this pointer.\par\medskip
\ifdefined\hypertarget\hypertarget{sourcepointer-2}{}\fi
\noindent\textit{Source reference.} Original reference 01ZM: focused-seeds/geometry/S9-morph-descend-cd/ref-tags.txt, source lines 4--4. Navigation pointer only: the counted proof cites this statement by \textbackslash{}ref as a standing input; the deduction the counted proof has to make is carried out inside the counted span, and no part of that deduction is delegated to this pointer.\par\medskip
\ifdefined\hypertarget\hypertarget{sourcepointer-3}{}\fi
\noindent\textit{Source reference.} Original reference 0GTB: focused-seeds/geometry/S9-morph-descend-cd/ref-tags.txt, source lines 5--5. Navigation pointer only: the counted proof cites this statement by \textbackslash{}ref as a standing input; the deduction the counted proof has to make is carried out inside the counted span, and no part of that deduction is delegated to this pointer.\par\medskip
\ifdefined\hypertarget\hypertarget{sourcepointer-4}{}\fi
\noindent\textit{Source reference.} Original reference 07RR: focused-seeds/geometry/S9-morph-descend-cd/ref-tags.txt, source lines 6--6. Navigation pointer only: the counted proof cites this statement by \textbackslash{}ref as a standing input; the deduction the counted proof has to make is carried out inside the counted span, and no part of that deduction is delegated to this pointer.\par\medskip
\ifdefined\hypertarget\hypertarget{sourcepointer-5}{}\fi
\noindent\textit{Source reference.} Original reference 01ZC: focused-seeds/geometry/S9-morph-descend-cd/ref-tags.txt, source lines 7--7. Navigation pointer only: the counted proof cites this statement by \textbackslash{}ref as a standing input; the deduction the counted proof has to make is carried out inside the counted span, and no part of that deduction is delegated to this pointer.\par\medskip
\begin{lemma}
\label{lemma-descend-cd}
Let $S = \lim_{\lambda \in \Lambda} S_\lambda$
be a limit of a directed system of schemes with affine transition morphisms.
Let $0 \in \Lambda$ and let $f_0 : X_0 \to Y_0$
be a morphism of schemes over $S_0$.
For $\lambda \geq 0$ let $f_\lambda : X_\lambda \to Y_\lambda$
be the base change of $f_0$ to $S_\lambda$ and
let $f : X \to Y$ be the base change of $f_0$ to $S$. If
\begin{enumerate}
\item $f$ is completely decomposed,
\item $Y_0$ is quasi-compact and quasi-separated, and
\item $f_0$ is locally of finite presentation,
\end{enumerate}
then there exists an $\lambda \geq 0$ such that $f_\lambda$
is completely decomposed.
\end{lemma}

\begin{proof}
Since $Y_0$ is quasi-compact and quasi-separated, the scheme $Y$,
which is affine over $Y_0$, is quasi-compact and quasi-separated.
Choose $n \geq 0$ and $Z_i \to Y$, $i = 1, \ldots, n$ as in
Lemma \ref{lemma-decompose}. Denote $\sigma_i : Z_i \to X$
morphisms over $Y$ which exist by our choice of $Z_i$.
After increasing $0 \in \Lambda$ we may assume there exist
morphisms $Z_{i, 0} \to Y_0$ of finite presentation
whose base changes to $S$ are the morphisms $Z_i \to Y$, see
Limits, Lemma \ref{limits-lemma-descend-finite-presentation}.
By Limits, Lemma \ref{limits-lemma-descend-immersion}
we may assume, after possibly increasing $0$, that $Z_{i, 0} \to Y_0$
is an immersion. Since $\coprod Z_i \to Y$ is surjective, we may assume,
after possibly increasing $0$, that
$\coprod Z_{i, 0} \to Y_0$ is surjective, see
Limits, Lemma \ref{limits-lemma-descend-surjective}. Observe that
$Z_i = \lim_{\lambda \geq 0} Z_{i, \lambda}$
where $Z_{i, \lambda} = Y_\lambda \times_{Y_0} Z_{i, 0}$.
Let us view the compositions
$$
Z_i \xrightarrow{\sigma_i} X \to X_0
$$
as morphisms over $Y_0$. Since $f_0$ is locally of finite
presentation by Limits, Proposition
\ref{limits-proposition-characterize-locally-finite-presentation}
we can find a $\lambda \geq 0$ such that there exist
morphisms $\sigma'_{i, \lambda} : Z_{i, \lambda} \to X_0$
over $Y_0$ whose precomposition with $Z_i \to Z_{i, \lambda}$
are the displayed arrows. Of course, then $\sigma'_{i, \lambda}$
determines a morphism $\sigma_{i, \lambda} : Z_{i, \lambda} \to
X_\lambda = X_0 \times_{Y_0} Y_\lambda$ over $Y_\lambda$.
Since $\coprod Z_{i, \lambda} \to Y_\lambda$ is surjective
we conclude that $X_\lambda \to Y_\lambda$ is completely decomposed.
\end{proof}

\end{document}
